% Nội dung phần 6: Phụ lục.

\section{Phụ lục}

\subsection{Hiệu năng phân tầng của mô hình được chọn}

Bảng~\ref{tab:subgroup} kiểm tra xem Logistic GLM có hoạt động đồng đều trên các
phân nhóm hay không, tính trên temporal test 01--12/08.

\begin{table}[H]
\centering
\caption{Hiệu năng phân tầng của Logistic GLM trên temporal test 01--12/08}
\label{tab:subgroup}
\small
\begin{tabular}{llrrrrrr}
\toprule
\textbf{Phân tầng} & \textbf{Nhóm} & \textbf{N} & \textbf{Prevalence}
& \textbf{ROC--AUC} & \textbf{AP} & \textbf{Log loss} & \textbf{Brier} \\
\midrule
\multirow{4}{*}{Category}
 & Trousers & 4.067 & 12,59\% & 0,638 & 0,200 & 0,365 & 0,107 \\
 & Skirts   & 2.769 & 15,57\% & 0,632 & 0,216 & 0,420 & 0,129 \\
 & Blouses  & 3.205 & 13,92\% & 0,688 & 0,252 & 0,380 & 0,114 \\
 & Sale     & 3.907 & 13,16\% & 0,693 & 0,252 & 0,364 & 0,108 \\
\midrule
\multirow{6}{*}{Order}
 & 1--2   & 3.437 & 18,10\% & 0,612 & 0,241 & 0,465 & 0,146 \\
 & 3--4   & 2.357 & 15,27\% & 0,616 & 0,233 & 0,418 & 0,127 \\
 & 5--7   & 2.395 & 13,70\% & 0,631 & 0,213 & 0,388 & 0,116 \\
 & 8--12  & 2.273 & 13,33\% & 0,664 & 0,261 & 0,370 & 0,110 \\
 & 13--20 & 1.653 & 9,13\%  & 0,670 & 0,198 & 0,292 & 0,080 \\
 & $>$20  & 1.833 & 7,58\%  & 0,723 & 0,187 & 0,245 & 0,066 \\
\bottomrule
\end{tabular}
\end{table}

Hiệu năng \emph{không} đồng đều giữa các phân nhóm. Theo danh mục, ROC--AUC dao
động từ 0,632 (Skirts) đến 0,693 (Sale). Theo tiến trình session, xu hướng rõ
ràng hơn: AUC tăng đều từ 0,612 ở khoảng \texttt{Order} 1--2 lên 0,723 với các
click sau vị trí 20. Nói cách khác, mô hình dự báo tốt hơn ở giai đoạn muộn của
phiên -- đúng lúc thông tin lịch sử duyệt đã tích lũy đủ. Đây cũng là hạn chế
thực tiễn: mô hình yếu nhất chính ở giai đoạn đầu, nơi rủi ro rời bỏ cao nhất.

Lưu ý rằng Brier score giảm dần theo \texttt{Order} chủ yếu vì prevalence giảm
(từ 18,10\% xuống 7,58\%), không phải vì mô hình tốt lên; hai chỉ số này phải đọc
cùng nhau.

\subsection{Phân công công việc}

\begin{table}[H]
\centering
\caption{Phân công công việc trong nhóm}
\label{tab:assignment}
\small
\begin{tabular}{>{\raggedright\arraybackslash}p{0.22\textwidth}
                >{\raggedright\arraybackslash}p{0.14\textwidth}
                >{\raggedright\arraybackslash}p{0.52\textwidth}}
\toprule
\textbf{Họ và tên} & \textbf{MSSV} & \textbf{Nội dung phụ trách} \\
\midrule
Nghi & 25C230XX &
Phần 1: giới thiệu bộ dữ liệu, xác định mục tiêu phân tích và xây dựng bản đề
xuất phương pháp; phối hợp viết Phần 4. \\[2pt]
Lê Hoàng Như & 25C23014 &
Phần 2: làm sạch dữ liệu, thống kê mô tả và trực quan hóa; phối hợp viết Phần 4;
soạn slide thuyết trình. \\[2pt]
Nguyễn Nhật Quang & 25C23002 &
Phần 3: xây dựng và đánh giá mô hình, toàn bộ bảng kết quả và biểu đồ; tổng hợp
nội dung hai file Rmd thành báo cáo nộp; phối hợp viết Phần 4. \\
\bottomrule
\end{tabular}
\end{table}

% TODO: cập nhật MSSV và họ tên đầy đủ của thành viên Nghi ở bảng trên
% và trong \studentname của file main.tex.

\subsection{Ghi chú về tính tái lập}

Toàn bộ kết quả trong báo cáo được sinh từ một lần render duy nhất của file
\texttt{de\_xuat\_phan\_tich\_clickstream.Rmd} bằng R~\cite{rcore2026} trong một
session sạch, không phụ thuộc vào object còn sót lại trong môi trường làm việc.

\begin{table}[H]
\centering
\caption{Môi trường thực thi}
\label{tab:environment}
\small
\begin{tabular}{ll}
\toprule
\textbf{Thành phần} & \textbf{Phiên bản} \\
\midrule
R           & 4.6.0 (2026-04-24) \\
rmarkdown   & 2.31 \\
knitr       & 1.51 \\
mgcv        & 1.9.4 \\
sandwich    & 3.1.1 \\
lmtest      & 0.9.40 \\
car         & 3.1.5 \\
dplyr       & 1.2.1 \\
tidyr       & 1.3.2 \\
ggplot2     & 4.0.3 \\
patchwork   & 1.3.2 \\
scales      & 1.4.0 \\
emmeans     & 2.0.4 \\
\bottomrule
\end{tabular}
\end{table}

\textbf{Nguồn dữ liệu.} File \texttt{e-shop clothing 2008.csv} (phân tách bằng
dấu chấm phẩy, 14 cột, 165.474 dòng), không chỉnh sửa thủ công.

\textbf{Seed ngẫu nhiên.} Mọi thủ tục có yếu tố ngẫu nhiên đều dùng
\texttt{set.seed(20260726)}: chia five-fold cross-validation theo session, 1.000
lần cluster bootstrap cho risk difference thô, và 500 lần cluster bootstrap cho
khoảng tin cậy của các metric trên temporal test. Nhờ vậy toàn bộ con số trong
báo cáo tái lập được chính xác.

\textbf{Kiểm tra tự động trong pipeline.} File Rmd chứa các khẳng định
\texttt{stopifnot()} được kiểm tra ở mỗi lần chạy: đúng 165.474 dòng và 24.026
session; không có giá trị khuyết thiếu hay dòng trùng lặp; \texttt{order} liên
tục trong mọi session; mỗi session có đúng một click \texttt{Exit = 1}; và không
có session nào xuất hiện đồng thời ở hai tập dữ liệu. Nếu bất kỳ điều kiện nào
sai, quá trình render sẽ dừng với lỗi thay vì tạo ra báo cáo không hợp lệ.

\textbf{Lệnh render.} Chạy từ thư mục chứa file Rmd và file CSV:

\begin{lstlisting}[language=R]
rmarkdown::render(
  "de_xuat_phan_tich_clickstream.Rmd",
  params = list(show_code = FALSE),
  envir  = new.env()
)
\end{lstlisting}

Tham số \texttt{show\_code = FALSE} tạo bản nộp không hiển thị mã nguồn; đặt
\texttt{show\_code = TRUE} sẽ tạo bản kiểm tra có kèm mã R cho từng bảng và biểu
đồ. Các hình trong báo cáo LaTeX này được xuất ở độ phân giải 200 dpi từ đúng lần
chạy nói trên và lưu trong thư mục \texttt{img/} với tiền tố \texttt{fig-}.
